%% The migration case: [lazy,langsci], one document, both syntaxes.
%%
%% This is what the option exists for -- a document written in the dot
%% syntax moving to \ea one example at a time -- so what it asserts is that
%% converting an example changes nothing around it: the counter runs on
%% through all three syntaxes, the geometry of a converted example is the
%% geometry of the one before it, and the cross-references still resolve.
%%
%% It also pins the two carve-outs from the one-syntax-per-example rule.
%% Within an exe batch, \a. ... \z. keeps working: that is a documented and
%% separately tested promise of [lazy,gb4e] (see the manual, "Mixing the two
%% syntaxes"), and [langsci] must not have narrowed it.  And a footnote is a
%% new stream, so a footnote hung on an example written one way may hold an
%% example written the other.
%%
%% (\ea in running text opens an example; the prose here avoids the name.)
\input{_preamble-langsci}
\begin{document}
\ex.\label{mx:dot} MXDOT an example still in the dot syntax.
\a. MXDOTA its first sub-example.
\b. MXDOTB its second.
\z.

\ea\label{mx:ea} MXEA the next example, already converted.
  \ea MXEAA its first sub-example.
  \ex MXEAB its second.
  \z
\z

\begin{exe}
\ex MXEXE an exe batch, which the ban does not touch.
\a. MXEXEA a dot sub-level inside the batch.
\z.
\ex MXEXEB back at the batch's main level.
\end{exe}

\ea MXFN an example carrying a footnote.\footnote{MXFNTEXT the note, which
holds an example of its own, written the other way.

\ex. MXFNEX a dot-syntax example inside the footnote.

}
\z

Refs: \ref{mx:dot}, \ref{mx:ea} and \Last.
\end{document}
